% ---
% id: meat-core/cor:constrained-key-length
% kind: cor
% title: "Key length from the constrained value"
% summary: >-
%   Under the attack-set licence the accumulated entropy is at least n times
%   the value of the program under the licensed list, less the
%   acceptance-probability penalty. That is the form the key-length theorem
%   reads, so its key length holds with the constrained value.
% prominence: standard
% status: implemented
% uses: [thm:constrained-accumulation, cor:reduction-constrained, thm:keylength]
% assumes: [A-ATK]
% refs: []
% provenance: newly-derived
% implements: certiqs_key.meat_core.finitekey.keylength:key_length_permitted
% lean:
% ---

\begin{corollary}
\label{cor:constrained-key-length}
Under the hypotheses of Theorem~\ref{thm:constrained-accumulation}, with
Corollary~\ref{cor:reduction-constrained},
\begin{equation*}
  \Hsandup\big(Z_1^n\big|\Chat_1^nI_1^n\Ehat\big)_{\rho_{|\OmAT}}
  \;\ge\;n\,\hlow(\CList)
  -\frac{\alpha}{\alpha-1}\log_2\frac{1}{\Pr[\OmAT]},
\end{equation*}
which is the form Theorem~\ref{thm:keylength} consumes; under
A-ATK (\xref{protocol-compiler/ass:atk}) and the remaining
hypotheses of Theorem~\ref{thm:keylength}, the key-length statement
of Theorem~\ref{thm:keylength} therefore holds with the constrained $\hlow$
(its constrained branch).
\end{corollary}

\begin{proof}
Theorem~\ref{thm:constrained-accumulation} gives
$\Hsandup(Z_1^n|\Chat_1^nI_1^n\Ehat)_{\rho_{|\OmAT}}
\ge n\,\hqkd(\CList)-\frac{\alpha}{\alpha-1}\log_2\frac{1}{\Pr[\OmAT]}$
under its hypotheses, and Corollary~\ref{cor:reduction-constrained} gives
$\hqkd(\CList)\ge\hlow(\CList)$ for the same licensed list; since $n\ge1$,
substituting the second inequality into the first yields the display. That
display is the accumulation bound the secrecy step of
Theorem~\ref{thm:keylength} reads, with $\hlow(\CList)$ in the place of
$\hlow$, so the constrained branch of that theorem applies under the
hypotheses of Theorem~\ref{thm:constrained-accumulation}, which include A-ATK,
together with the remaining hypotheses of Theorem~\ref{thm:keylength}: its
post-processing model and the assumptions it names, $\alpha\in(1,2)$ and
$\epa,\eev\in(0,1]$.
\end{proof}
